Pomocou 47 nukleidov dlhy úvodný úsek som prehlasil ako primer a rozhodol som sa ho vyhladať v každom čítaní, nielen na začiatku sekvencie, ale aj v jej ďalších častiach, aby som zachytil aj prípady jeho posunu alebo výskytu vo vnútri sekvencie.

subor obsahujuci citania z dorada v formate fasta som načítal do pamäti a následne som použil algoritmus z bioAlignments knižnice Biopython, ktorý umožňuje vyhľadávať lokálne zarovnania medzi primerom a jednotlivými čítaniami.

aligner.align(primer, read) vracia všetky lokálne zarovnania medzi primerom a daným čítaním zoradené podľa skóre zarovnania.
kazdy alignment obsahuje informácie o pozícii zarovnania v celkovom reťazci a lokalne zarovnanom retazci, ktorý zodpovedá v nasom prípade primeru.
je zapisaný v 3 dimenzionálnom formáte, 2 N 2, kde 1 riadok obsahuje lokalne zarovnany retazec, 2 riadok obsahuje celkovy retazec. 
potom kazde pole obsahuje zarvnaie odkial pokial je cast zarovnana a je to rozdelene na N fragmentov, kde N je počet zarovnaných úsekov ten sa zhoduje s počtom fragmentov v lokalne zarovnanom retazci.

Nato aby sme zistili odkial pokial sa lokalne zarovnaný retazec zhoduje s primerom tak data najdeme tu aligned[1][0][0] a aligned[1][-1][1]
/input{chapters/CodeSnipets/snippets/python-find-occurrences.tex}

primery ulozim do formatu [[], [65], [55], [], [64], [], [55], [74, 73], [64], [75], [72, 70], [], [], [75], [39], [75], [], [], [72].. 
kde kazdy element zodpoveda jednému čítaniu a obsahuje pozície, kde sa primer nachádza v danom čítaní. Ako je vidno nie v kazdom retazci sa primer nachadza, ale v niektorych sa nachadza viac krat.
je to zdvovodu bud ostatne citania obsahovali komplementarnu sekvenciu alebo sa precitali v opacnom smere, alebo sa primer nenachadza zdvovodu pretrhnutia dna retazca v roznych miestach 
alebo sa nachadza ale je prilis odlišný ze sa dostal pod mojim nastaveným thresholdom pre zarovnanie teda ho ignorujem.

\begin{lstlisting}[language=Python, caption={Stripping primers from Clover input}, label={lst:python-strip-primers}]
    with open(clover_input_path, "r") as f:
        lines = f.readlines()

    with open(PrimersPath, "r") as f:
        datas = f.read()

    primers = ast.literal_eval(datas)

    newLines = []
    for (line, linePrimes) in zip(lines,primers):
        prevValue = -1
        for prime in linePrimes:
            if(prevValue!=-1 ):
                if prevValue-prime<50:
                    continue
            newline = line[prime:prevValue]
            if(len(newline)<50):
                continue
            newLines.append(newline)
            prevValue = prime

    WriteLines=[]
    for (i,line) in enumerate(newLines):
        WriteLines.append(f"{i} {line}\n")

    with open("PrimersStripClover.txt", "w") as f:
        f.writelines(WriteLines)
\end{lstlisting}
Tak som zobral useky za primerom a zapisal ich do noveho suboru, ktory obsahuje citania od primeru az po koniec retazca. Alebod do dalsieho primeru ak sa nachadzal viac krat.
Tieto data som následne znovu zoskupil pomocou algoritmu Clover, ale tentokrát som použil iba zvyšné časti sekvencií bez DNA primerov.
Co sposobilo ze tentokrát klastre obsahovali sekvencie, ktoré mali zhodu nielen na začiatku, ale aj ďalej v reťazci, co indikovalo ze sa sekvencie pochadzaju z nasledkom polimerizacie z toho isteho dna retazca v spravnom smere a nie z toho isteho dna retazca v opacnom smere alebo z komplementarneho retazca.


Na Klastre som aplikoval viacnásobné zarovnanie.
\begin{lstlisting}[language=Python, caption={Clustering without primers}, label={lst:python-clustering-without-primers}]
ClusterFile = TextFileM(path="clusters_output.txt")
ClusterFile = ClusterFile.load()
ClusterFile = parse_tuple_file(ClusterFile)
ClusterFile = get_clusters(ClusterFile)
clusters = ClusterFile.data

#ClusterFile = pick_largest_cluster(ClusterFile)
#LargestCluster = []
#LargestCluster = ClusterFile.data

clusters = [l for l in clusters if len(l) > 10]
count=1
for clust in clusters:
    ClusterFile.data=clust
    ClusterFile = retrieve_dna_lines(ClusterFile, clover_input_path)
    DataToBeAligned = ClusterFile.data

    ClusterSize = len(DataToBeAligned)
    (Aligned, Unaligned) = DataBeamAlgorithmSplit(DataToBeAligned, 10)
    dist = len(Aligned[0].replace("_",""))
    if len(Aligned[0].replace("_",""))<10:
        continue

    Aligned.insert(0,str(clust))

    TextFileM(data=Aligned).write_data_to_file(f"GeneratedData/Aligned_dnaSize-{dist}num-{ClusterSize}id-{count}.txt")
    count+=1
\end{lstlisting}
Ukázalo sa, že klastre obsahovali sekvencie, ktoré mali zhodu nielen na začiatku, ale aj ďalej v reťazci.
Niektoré klastre obsahovali sekvencie, ktoré sa zhodovali až do dĺžky 100 nukleotidov, čo je výrazně ďalší úsek než pôvodní primer.


%